% !TeX root = ../../../main.tex
이번 학기에 무려 23학점을 듣고 있는 기룡이는 매일 아침 지하철을 타고 1교시를 들으러 학교에 가야 한다.
하지만 학교에 가는 게 등교인지 등산인지 구분하지 못할 정도로 힘들었던 기룡이는 이런 상상을 해보았다.

``지하철역과 학교의 모든 건물의 지하에 거대한 도시를 만든다면 어떨까?
엘리베이터를 타고 내려가서 걷기만 하면 되니까 엄청 편할 텐데!
음... 걷는 것도 힘드니까 가장 먼 건물 사이에는 무빙워크까지 설치하면 좋겠다!
지하 도시를 만드는 건 돈이 엄청 많이 들 텐데, 아무래도 지하 도시는 넓이가 최소인 볼록 다각형이 돼야 할 것 같아.
그렇다면 이 도시의 넓이와 무빙워크의 길이는 얼마가 될까?''

지하철역을 포함한 학교의 모든 건물은 2차원 평면 위의 점 $(x, y)$로 표현되며, 두 점 $u = (x_u, y_u)$, $v = (x_v, y_v)$ 사이의 거리는 $\sqrt{(x_u-x_v)^2+(y_u-y_v)^2}$로 정의된다.

지하 도시는 모든 건물을 포함하는 볼록 다각형 중 넓이가 최소인 것으로 생각할 수 있으며, 무빙워크의 길이는 주어진 건물들 사이의 최장 거리로 정의된다.
단일 점의 경우, 넓이와 길이가 모두 0인 볼록 다각형이며, 단일 선분의 경우 넓이가 0인 볼록 다각형으로 본다.

기룡이를 위해 도시의 넓이와 무빙워크의 길이를 구해보자.
